\selectlanguage{french}
\settitle[AndroLog : prérequis]{AndroLog : prérequis pour son utilisation}
\setauthor[J. Samhi]{Jordan Samhi\\\email{jordan.samhi@uni.lu}}
\institute{Université du Luxembourg}
\maketitle

AndroLog ne nécessite aucun logiciel payant ni matériel spécialisé.

\section{Instrumenter et analyser une app}

Pour instrumenter ou analyser un APK, AndroLog nécessite le \textbf{SDK Android}, avec une plateforme Android. L’instrumentation nécessite également les \textit{Build Tools}, notamment \texttt{zipalign} et \texttt{apksigner}. Ces outils sont disponibles gratuitement ; leur utilisation est soumise aux conditions de licence du SDK Android.

AndroLog nécessite aussi \textbf{Java 17 ou une version plus récente}.

\section{Exécuter l’app et collecter les logs}

Pour exécuter une app instrumentée et collecter de nouveaux logs, il faut un \textbf{émulateur Android ou un téléphone Android}, ainsi que l’outil \texttt{adb} fourni dans les \textit{Platform Tools} du SDK. Sur un téléphone physique, le débogage USB doit être activé et l’ordinateur autorisé. Les images d’émulateur sont soumises à leurs propres conditions de licence.

\section{Plateformes et consultation d’un rapport}

AndroLog est multiplateforme. Les sources permettent de construire l’outil pour macOS, Windows et Linux ; la compilation nécessite notamment un JDK, Maven, Node.js et npm.

Un rapport HTML déjà généré s’ouvre dans un navigateur web, sans SDK Android, appareil Android ou environnement Java.
